\section{Summary of Results}

\subsection{Interpretation of Results}

The voxel-sampler-based pipeline more than halved mean elapsed processing time in the DALES benchmarks. Its lower overall mIoU was accompanied by different outcomes across classes: buildings favored the proposed workflow, while power lines favored probability sampling (Tables~\ref{tab:dales_probability} and~\ref{tab:dales_voxel}). The practical trade-off therefore depends on which objects the application needs to recover, as well as how quickly a~scene must be processed.

On BRIK, the proposed workflow achieved higher mIoU despite lower overall accuracy. Its largest F1 gains concerned poles and support devices, both relevant to the railway-corridor application for which it was developed. Railway tracks favored probability sampling, while overhead-line F1 was similar for both workflows (Tables~\ref{tab:brik_voxel} and~\ref{tab:brik_probability}). These differences make per-class results particularly relevant on the imbalanced BRIK data: a~single accuracy value gives little indication of how well individual infrastructure classes are recovered. The opposite mIoU rankings on BRIK and DALES also suggest that the choice of sampling workflow should account for the scene characteristics.

The density-reduction results described in the Results section also favor the proposed workflow: across the tested voxel sizes, its accuracy declined less and its returned-prediction coverage remained complete. Fixed spatial windows may help retain a~more consistent physical context as density changes, whereas fixed-cardinality neighborhoods expand in sparse regions. This is a~possible explanation for the reported trend; the accuracy values refer to the labeled points with returned predictions, whose sets vary between configurations.

\subsection{Advantages and Limitations of the Solution}

The largest timing difference occurs in model prediction, which takes only a~small part of the proposed pipeline's total elapsed time. A~finite grid of windows provides an explicit spatial processing schedule without a~coverage-driven query loop. Together with the class-level gains on BRIK and for DALES buildings, this makes the proposed workflow useful where processing time and the recovery of these classes are priorities. Its complete returned-prediction coverage in the density-reduction experiment is a~further practical advantage, although coverage describes label availability rather than label correctness.

The shorter elapsed time comes with higher memory demand, as reflected in peak PSS. The much larger process-tree RSS can count shared pages repeatedly and should be distinguished from physical memory use. The reduction in elapsed time was much larger than the decrease in accumulated CPU time. With model prediction occupying a~small part of the elapsed time, preparation, aggregation, and reconstruction become relevant optimization targets.

Fixed spatial extent introduces its own constraints. Windows below the point-count threshold bypass inference, and reconstruction transfers scores from nearby retained points. Sparse structures and class boundaries may therefore be sensitive to the window size, sampling threshold, and interpolation scale. Outer XY tiles have no overlap margin, which may also limit context near their boundaries. The lower DALES power-line recall and BRIK railway-track recall motivate examining these mechanisms, while the low truck recall in both pipelines identifies a~difficulty shared by the compared workflows.

\subsection{Possible Extensions}

Further development could focus on preparation and reconstruction, where improvements would complement the already short model prediction stage. Spatial indexing could replace repeated point-in-window checks, and sharing a~spatial index between reconstruction workers, where supported, could reduce duplicated indexing work and storage. These changes should be evaluated using total elapsed time and PSS, alongside model prediction time.

A~second direction is to adjust window extent, overlap, and the minimum point count jointly with the reconstruction neighborhood and distance scale. Tests focused on power lines, railway tracks, and sparse windows could examine whether additional context or less spatial smoothing improves recall while retaining the observed gains for buildings and railway infrastructure. Adding margins to outer tiles is another candidate for testing boundary effects. Future comparisons should report per-class quality together with prediction coverage on a~common evaluation set, so that parameter choices reflect both the objects recovered and the processing cost.